% =====================================================================
%  Atividade 4 - Métodos (seminário de dissertação) - versão didática
%  Notação alinhada às Aulas 6, 7 e 8 de Microeconometria II (Profa. Dani, PPGE-UFF)
%  Compilar no Overleaf com pdfLaTeX (padrão). Limite: 3 páginas de texto.
% =====================================================================
\documentclass[12pt,a4paper]{article}

\usepackage[utf8]{inputenc}
\usepackage[T1]{fontenc}
\usepackage[brazil]{babel}
\addto\captionsbrazil{\renewcommand{\tablename}{Quadro}}

\usepackage{mathptmx}
\usepackage{amsmath}
\usepackage[top=3cm, left=3cm, bottom=2cm, right=2cm]{geometry}
\usepackage{setspace}
\setstretch{1.12}
\usepackage{indentfirst}
\setlength{\parindent}{1.25cm}
\usepackage{microtype}

\usepackage{booktabs}
\usepackage{tabularx}
\usepackage{float}
\usepackage{titlesec}
\titleformat{\section}{\normalfont\bfseries}{\thesection}{0.5em}{}
\titlespacing*{\section}{0pt}{8pt}{3pt}
\usepackage[hidelinks]{hyperref}
\newcolumntype{L}{>{\raggedright\arraybackslash}X}

\newenvironment{referencias}{%
  \section*{Referências}
  \begin{singlespace}\footnotesize
  \setlength{\parindent}{0pt}
  \setlength{\parskip}{2pt}
  \raggedright
}{%
  \end{singlespace}
}

\begin{document}

\begin{center}
  {\bfseries ATIVIDADE 4 -- MÉTODOS\\
  O EFEITO DA COP30 SOBRE O EMPREGO FORMAL EM ALOJAMENTO E ALIMENTAÇÃO}\\[4pt]
  Marcela Paz da Silva Enriquez
\end{center}

\section{A pergunta e o problema de identificação}

A pesquisa quer saber se a COP30, realizada em Belém em novembro de 2025, mudou o emprego formal
em Alojamento e Alimentação (hotéis, restaurantes e bares; divisões 55 e 56 da CNAE) e se esse efeito
permaneceu depois que o evento acabou. Para responder, seria preciso comparar dois mundos: a Belém que
sediou a COP30 e a mesma Belém, no mesmo mês, sem a COP30. Em notação, $y_{it}(1)$ é o emprego da
capital $i$ no mês $t$ com o evento e $y_{it}(0)$ é o emprego sem o evento; a variável $d_i$ indica
quem recebeu o tratamento ($d_i = 1$ para Belém e $d_i = 0$ para as demais capitais). O efeito da COP30
em Belém é a diferença entre os dois mundos, $y_{it}(1) - y_{it}(0)$, e a sua média,
$ATT(t) = E[y_{it}(1) - y_{it}(0) \mid d_i = 1]$, é o efeito médio do tratamento sobre os tratados: o
efeito do evento na cidade que de fato o recebeu. O problema é que só um dos mundos aconteceu. O
emprego com a COP30 é observado nos dados, mas o emprego que Belém teria sem ela, $y_{it}(0)$, nunca
pode ser visto. Esse é o problema fundamental da inferência causal, e o método serve para estimar esse
cenário que falta, chamado de contrafactual.

Comparar Belém com as outras capitais em um único mês daria o ATT mais um viés de seleção, porque
Belém não foi escolhida ao acaso. Métodos que controlam apenas características observáveis, como o
pareamento, não resolvem o problema, pois a escolha da sede depende de fatores não observados.
Também não há um limiar que defina o tratamento, como
exige a regressão descontínua. Como há dados de muitos meses, a estratégia adequada é a de \textbf{diferenças em diferenças com dados em painel}, que elimina as características
não observadas fixas no tempo.

\section{Dados}

Os dados vêm dos microdados do Novo CAGED, do Ministério do Trabalho e Emprego, que registram cada
contratação e demissão com carteira assinada, com município e atividade da empresa. Para cada capital
e mês calcula-se o saldo (contratações menos demissões) em Alojamento e Alimentação, dividido pela média
mensal de contratações da própria capital antes do evento; assim, $y_{it} = 10$ significa ``10\% de um
mês normal de contratações''. A amostra reúne as capitais brasileiras: Belém
e 25 capitais de controle ($n = 26$). O Rio de Janeiro fica fora porque sediou a Cúpula do G20 em
novembro de 2024, um evento do mesmo tipo; como também recebeu um choque de demanda, não representa
o cenário ``sem evento'' que o grupo de controle precisa representar. As mesmas 26 capitais são
observadas em todos os 55 meses (jan/2022 a jul/2026), sem falhas: um \textbf{painel balanceado}.

\section{Método: diferenças em diferenças}

\textbf{O modelo canônico.} Seja $T_t = 1$ nos meses depois do início do tratamento e $T_t = 0$ antes.
A regressão de diferenças em diferenças é
\begin{equation}
y_{it} = \alpha + \theta d_i + \lambda T_t + \gamma\,(d_i \times T_t) + \epsilon_{it},
\end{equation}
e o Quadro 1 mostra o significado de cada coeficiente.

\begin{table}[H]
\centering\small
\caption{Interpretação dos coeficientes}
\begin{tabularx}{\linewidth}{lccL}
\toprule
 & \textbf{Antes ($T_t = 0$)} & \textbf{Depois ($T_t = 1$)} & \textbf{Depois menos antes}\\
\midrule
Belém ($d_i = 1$) & $\alpha + \theta$ & $\alpha + \theta + \lambda + \gamma$ & $\lambda + \gamma$\\
Outras capitais ($d_i = 0$) & $\alpha$ & $\alpha + \lambda$ & $\lambda$\\
\midrule
\textbf{Diferença em diferenças} & & & $\gamma^{DID} = \Delta Tr - \Delta Cr = \boldsymbol{\gamma}$\\
\bottomrule
\end{tabularx}
\end{table}

O coeficiente $\theta$ é a diferença permanente entre Belém e as outras capitais, e some na diferença
``depois menos antes''; $\lambda$ é o choque agregado que atinge todas as capitais, e some na
diferença ``Belém menos as outras''. Sobra $\gamma$, o efeito da COP30.

\textbf{Vários meses: efeitos fixos e $ATT(t)$.} Como o painel tem muitos meses antes e depois, o modelo
é estimado em efeitos fixos em dois sentidos (TWFE), com uma interação para cada mês pós-tratamento:
\begin{equation}
y_{it} = \alpha_i + \lambda_t + \sum_{\tau} \gamma_\tau \, \mathbf{1}_\tau(d_i) + \epsilon_{it}.
\end{equation}
O efeito fixo $\alpha_i$ generaliza $\theta$: capta as características não observadas e constantes no
tempo de cada capital, inclusive as que levaram Belém a ser escolhida. Ele é definido por capital e mês
do ano, para que cada cidade tenha a sua própria sazonalidade (em Belém, o Círio de Nazaré em outubro).
O efeito fixo $\lambda_t$ generaliza $\lambda$: capta choques agregados como juros, inflação e o ciclo
nacional. Cada $\gamma_\tau$ estima $ATT(\tau)$, o efeito em um mês específico, e a trajetória mostra se
o efeito cresce, se mantém ou se dissipa. Os efeitos são resumidos em três fases, preparação
(jul/2024--out/2025), evento (nov/2025) e desmobilização (dez/2025--jul/2026), pela média dos
$ATT(t)$ de cada fase, $\Gamma = \sum_t w(t)\,ATT(t)$, com pesos iguais dentro da fase. Como só Belém é tratada, e numa única
data, o grupo de controle é formado apenas por cidades que nunca sediaram a COP; por isso não ocorre o
viés do TWFE apontado por Goodman-Bacon (2021) para tratamentos que começam em datas diferentes.

\textbf{Hipóteses.} (1) \textit{Tendências paralelas}: sem a COP30, a evolução de $y_{it}(0)$ em
Belém seria igual à das outras capitais, em todos os meses; os níveis podem ser diferentes. Com vários
meses pré-tratamento, isso é verificado estimando $\gamma_\tau$ também antes de jul/2024, que devem ser
próximos de zero. (2) \textit{Não antecipação}: como o evento foi anunciado com mais de um ano de
antecedência, o tratamento começa no início da preparação (jul/2024), e não em novembro de 2025.
(3) \textit{SUTVA}: o emprego das capitais de controle não pode ter sido afetado pela COP30; por isso o
Rio é excluído e as capitais do Norte são testadas fora da amostra. (4) \textit{Irreversibilidade}:
Belém permanece como ``cidade que sediou a COP'' depois do evento, e os $ATT(t)$ da desmobilização medem
quanto do efeito sobrevive.

\textbf{Em que artigos o modelo se baseia.} O modelo segue diretamente Leoni e Boto-García (2023), que
avaliam a erupção vulcânica de La Palma, um choque temporário sobre uma economia dependente do turismo,
com um painel mensal de zonas turísticas das Ilhas Canárias. Como aqui, há uma única unidade tratada
(La Palma), efeitos fixos de unidade e de mês, e o tratamento é dividido em dois períodos, durante e
depois do evento, o que permite ver se o efeito persiste. Os autores encontram queda de cerca de
30\% no emprego da hotelaria durante e depois da erupção. A lógica geral é a de Card e Krueger (1994); no Brasil, o
trabalho mais próximo é Mancha (2026), apresentado no 48º Encontro Brasileiro de Econometria.

\textbf{Análise complementar: transbordamento no Pará.} Leoni e Boto-García (2023) acrescentam um grupo
``potencialmente tratado'' (as ilhas vizinhas) para medir se o efeito se espalhou. A mesma ideia é
aplicada aos municípios do Pará: Belém é a tratada, os demais municípios da Região Metropolitana são os
potencialmente tratados e o interior é o controle. Um coeficiente positivo para a Região Metropolitana
indica que o emprego transbordou para os vizinhos; um negativo, que Belém atraiu trabalhadores deles.

\section{Cuidados na estimação}

\begin{table}[H]
\centering\small
\caption{Cuidados do modelo de diferenças em diferenças e como são tratados}
\begin{tabularx}{\linewidth}{p{4.2cm}L}
\toprule
\textbf{Cuidado} & \textbf{Como é tratado nesta pesquisa}\\
\midrule
Tendência comum & Placebo de \textit{timing} (``evento'' falso em nov/2023, sem efeito esperado) e placebo de grupo (cada capital tratada como se tivesse sediado a COP30)\\
Transformações não lineares & O saldo pode ser negativo, por isso é usado em nível normalizado, e não em logaritmo\\
Composição dos grupos & Painel balanceado, com as mesmas 26 capitais em todos os meses\\
Inferência e correlação serial (Bertrand, Duflo e Mullainathan, 2004) & Com uma única unidade tratada, o p-valor vem do placebo de grupo: a posição de Belém no ranking dos 26 efeitos (Conley e Taber, 2011)\\
Outros choques em Belém & Outubro a dezembro estimados como um bloco (Círio); retirada de nov/2024 (Conferência Amazônia e Novas Economias)\\
\bottomrule
\end{tabularx}
\end{table}

Como há uma só cidade tratada, o resultado será conferido por um segundo método, o controle sintético
(Abadie, Diamond e Hainmueller, 2010). Em vez da média das outras capitais, ele usa como comparação uma
``Belém sintética'': uma média ponderada de capitais escolhida para ter o emprego mais parecido possível
com o de Belém antes da COP30. Se os dois métodos apontarem na mesma direção, a conclusão é mais robusta.

\begin{referencias}
ABADIE, A.; DIAMOND, A.; HAINMUELLER, J. Synthetic control methods for comparative case studies.
\textit{Journal of the American Statistical Association}, v. 105, n. 490, p. 493--505, 2010.

BERTRAND, M.; DUFLO, E.; MULLAINATHAN, S. How much should we trust differences-in-differences
estimates? \textit{Quarterly Journal of Economics}, v. 119, n. 1, p. 249--275, 2004.

CARD, D.; KRUEGER, A. B. Minimum wages and employment: a case study of the fast-food industry in New
Jersey and Pennsylvania. \textit{American Economic Review}, v. 84, n. 4, p. 772--793, 1994.

CONLEY, T. G.; TABER, C. R. Inference with ``difference in differences'' with a small number of
policy changes. \textit{Review of Economics and Statistics}, v. 93, n. 1, p. 113--125, 2011.

GOODMAN-BACON, A. Difference-in-differences with variation in treatment timing. \textit{Journal of
Econometrics}, v. 225, n. 2, p. 254--277, 2021.

LEONI, V.; BOTO-GARCÍA, D. The effect of natural disasters on hotel demand, supply and labour
markets: evidence from the La Palma volcano eruption. \textit{Environmental and Resource Economics},
v. 86, n. 4, p. 755--780, 2023.

MANCHA, A. The State strikes back: evidence from police crackdowns in São Paulo, Brazil. In: ENCONTRO
BRASILEIRO DE ECONOMETRIA, 48., 2026. \textit{Anais}. [S. l.]: SBE, 2026.
\end{referencias}

\end{document}
